\subsection{DEFINITION OF NORMAL SPACES}

Axiom \((\mathbf{O}_{\mathbf{IV}})\) for uniformizable spaces (§ 1, no. 5) can be stated in the following form: given any closed set A and any point \(x \in \complement \mathbf{A}\) , there is a continuous mapping of \(\mathbf{X}\) into {[}0, 1{]} which is equal to o at x and is equal to i at every point of A; this property can again be expressed by saying that in a uniformizable space we can separate a point and a closed set (not containing the point) by a continuous real-valued function.

We shall now study spaces in which it is possible in the same way to separate two disjoint closed sets by a continuous real-valued function:

DEFINITION 1. A topological space X is said to be normal if it is Hausdorff and satisfies the following axiom:

\((\mathbf{O}_{\mathbf{v}})\) If A and B are any two disjoint closed subsets of X, there exists a continuous mapping of X into {[}o, i{]} which is equal to o at every point of A and to i at every point of B.

Clearly every normal space is completely regular; but there are completely regular spaces which are not normal (see Exercises 9, 10, 13, 26 and § 5, Exercises 15 and 16).

The statement of Axiom \((\mathrm{O}_{\mathrm{V}})\) , like that of \((\mathrm{O}_{\mathrm{IV}})\) , involves the real line R as an auxiliary set. But there is a condition equivalent to \((\mathrm{O}_{\mathrm{V}})\) which involves no auxiliary set:

THEOREM 1 (Urysohn). Axiom \((\mathrm{O}_{\mathrm{V}})\) is equivalent to the following:

\((\mathbf{O}_{\mathbf{V}}^{\prime})\) If \(\mathbf{A}\) and \(\mathbf{B}\) are any two disjoint closed subsets of \(\mathbf{X}\) , then there exist two disjoint open sets \(\mathbf{U}, \mathbf{V}\) such that \(\mathbf{A} \subset \mathbf{U}\) and \(\mathbf{B} \subset \mathbf{V}\) .

It follows immediately that \((\mathrm{O}_{\mathrm{V}})\) implies \((\mathrm{O}_{\mathrm{V}}^{\prime})\) , for if f is a continuous mapping of X into {[}0, I{]} which is equal to o on A and to I on B, then the open sets \(\overline{f}^{1}([0, I/2[)\text{ and }\overline{f}^{1}([I/2,I])\text{ contain A and B respectively and do not intersect.}\)

To prove the converse, notice first that \((\mathrm{O}_{\mathrm{V}}^{\prime})\) is equivalent to the following axiom:

\((\mathbf{O}_{\mathbf{V}}^{\prime \prime})\) Given any closed set \(\mathbf{A}\) and any open neighbourhood \(\mathbf{V}\) of \(\mathbf{A}\) , there exists an open neighbourhood \(\mathbf{W}\) of \(\mathbf{A}\) such that \(\overline{\mathbf{W}} \subset \mathbf{V}\) .

If there is a continuous mapping \(f: \mathbf{X} \to [-\mathbf{i}, +\mathbf{i}]\) which is equal to \(-\mathbf{i}\) on \(\mathbf{A}\) and to \(+\mathbf{i}\) on \(\mathbf{B}\) , and if we put \(\mathbf{U}(t) = \overline{f}^1([-\mathbf{i}, t[)\) for each \(t \in [\mathbf{o}, \mathbf{i}]\) , then we have defined a family of open sets in \(\mathbf{X}\) , indexed by \([\mathbf{o}, \mathbf{i}]\) , such that (i) \(\mathbf{A} \subset \mathbf{U}(\mathbf{o})\) , (ii) \(\mathbf{B} \subset \mathbf{C}\mathbf{U}(\mathbf{i})\) and (iii) for each pair of real numbers \(t, t'\) such that \(\mathbf{o} \leqslant t < t' \leqslant \mathbf{i}\) we have

\[
\overline {{{\mathbf {U}}}} (t) \subset \mathbf {U} (t ^ {\prime});\tag{1}
\]

for \(\mathbf{U}(t)\) is contained in the closed set \(\overline{f}([-\mathrm{i}, t])\) . Conversely, suppose that we have defined a family \((\mathbf{U}(t))\) of open sets \((0 \leqslant t \leqslant 1)\) with these three properties (i), (ii) and (iii). For each \(x \in \mathbf{X}\) , put \(g(x) = 1\) if \(x \in \mathbb{C}\mathrm{U}(\mathrm{I})\) , and if \(x \in \mathrm{U}(\mathrm{I})\) let \(g(x)\) be the greatest lower bound of the values of \(t\) such that \(x \in \mathrm{U}(t)\) . Clearly \(0 \leqslant g(x) \leqslant 1\) for each \(x \in \mathrm{X}\) , \(g(x) = 0\) on \(\mathrm{A}\) , \(g(x) = 1\) on \(\mathrm{B}\) ; also \(g\) is continuous on \(\mathrm{X}\) , for if we put \(g(x) = a\) , we have \(|g(y) - g(x)| \leqslant \varepsilon\) for all \(y \in \mathrm{U}(a + \varepsilon) \cap \mathbb{C}\overline{\mathrm{U}}(a - \varepsilon)\) , which is a neighbourhood of \(x\) by (I) {[}with the conventions that \(\mathrm{U}(a + \varepsilon) = \mathrm{X}\) if \(a + \varepsilon > 1\) , and \(\mathrm{U}(a - \varepsilon) = \emptyset\) if \(a - \varepsilon < 0\) {]}.

Hence Theorem 1 will be proved if we can define a family \((\overline{\mathbf{U}}(t))\) of open sets satisfying conditions (i), (ii) and (iii) above; to do this we use Axiom \((\mathbf{O}_{\mathbf{V}}^{\prime \prime})\) . Take \(\mathbf{U}(\mathbf{i}) = \mathbb{C}\mathbf{B}\) ; since \(\mathbf{A} \subset \mathbf{U}(\mathbf{i})\) there exists an open set \(\mathbf{U}(\mathbf{o})\) such that \(\mathbf{A} \subset \mathbf{U}(\mathbf{o})\) and \(\mathbf{U}(\mathbf{o}) \subset \mathbf{U}(\mathbf{i})\) by \((\mathbf{O}_{\mathbf{V}}^{\prime \prime})\) . Suppose then that for each dyadic number \(k / 2^n\) ( \(k = \mathbf{o}, \mathbf{i}, \ldots, 2^n\) ) we have defined an open set \(\mathbf{U}(k / 2^n)\) , these sets being such that \(\overline{\mathbf{U}}(k / 2^n) \subset \mathbf{U}((k + \mathbf{i}) / 2)^n\) for \(\mathbf{o} \leqslant k \leqslant 2^n - \mathbf{i}\) . For each dyadic number

\[
(2 k + \mathrm{I}) / 2 ^ {n + 1} \quad (\mathrm{o} \leqslant k \leqslant 2 ^ {n} - \mathrm{I})
\]

there exists by \((\mathbf{O}_{\mathbf{v}}^{\prime \prime})\) an open set \(\mathrm{U}((2k + 1) / 2^{n + 1})\) such that

\[
\overline {{{{\mathbf {U}}}}} (k / 2 ^ {n}) \subset \mathrm{U} ((2 k + \mathrm{I}) / 2 ^ {n + 1})
\]

and

\[
\overline {{{\mathbf {U}}}} ((2 k + \mathrm{I}) / 2 ^ {n + 1}) \subset \mathbf {U} ((k + \mathrm{I}) / 2 ^ {n}).
\]

Hence for each dyadic number \(r\) such that \(0 \leqslant r \leqslant 1\) we can define an open set \(U(r)\) , such that \(A \subset U(o)\) , \(B \subset \complement U(1)\) , and

\[
\overline {{{\mathbf {U}}}} (r) \subset \mathbf {U} (r ^ {\prime})\tag{2}
\]

for each pair of dyadic numbers \(r, r'\) such that \(0 \leqslant r < r' \leqslant 1\) .

Now define, for each real number \(t \in [0, 1]\) ,

\[
\mathrm{U} (t) = \bigcup_ {r \leqslant t} \mathrm{U} (r) \quad (r \text { dyadic });
\]

by (2), this definition agrees with the preceding one for \(t\) dyadic; also, if \(0 \leqslant t < t' \leqslant 1\) , then there exist two dyadic numbers \(r, r'\) such that \(t \leqslant r < r' \leqslant t'\) ; by (2) we have \(\mathrm{U}(t) \subset \mathrm{U}(r)\) , hence

\[
\overline {{{\mathbf {U}}}} (t) \subset \overline {{{\mathbf {U}}}} (r) \subset \mathbf {U} (r ^ {\prime}) \subset \mathbf {U} (t ^ {\prime});
\]

this proves (1) and therefore completes the proof.

Theorem 1 will enable us to show that two important categories of topological spaces are normal. In the first place:

PROPOSITION 1. A compact space is normal.

A compact space satisfies axiom \((\mathrm{O}_{\mathrm{V}}^{\prime})\) , by Proposition 2 of Chapter I, § 9, no. 2.

As to locally compact spaces, every point of such a space has a compact neighbourhood, which is a normal subspace; but there are examples of locally compact spaces which are not normal (cf.~Exercises 9 and 26, and § 5, Exercise 15).

PROPOSITION 2. A metrizable space is normal.

Let X be a metrizable space and let d be a metric compatible with the topology of X. Let A, B be two disjoint closed subsets of X; since the functions \(d(x, \mathbf{A})\) and \(d(x, \mathbf{B})\) are continuous, the set U (resp. V) of points x such that \(d(x, \mathbf{A}) < d(x, \mathbf{B})\) {[}resp. \(d(x, \mathbf{B}) < d(x, \mathbf{A})\) {]} is open; clearly \(A \subset U\) and \(B \subset V\) and \(U \cap V = \varnothing\) , hence Axiom \((\mathrm{O}_{\mathrm{V}}^{\prime})\) is satisfied.

Remarks. 1) Proposition 2 gives another necessary condition for metrizability; but this condition, even in conjunction with all the necessary conditions given in § 2, does not give a set of sufficient conditions for metrizability (cf.~Exercise 6 and § 5, Exercise 10).

\begin{enumerate}
\def\labelenumi{\arabic{enumi})}
\setcounter{enumi}{1}
\tightlist
\item
  There are examples of normal spaces which are neither metrizable nor locally compact (see § 5, Exercise 16).
\end{enumerate}

By \((\mathrm{O}_{\mathrm{V}}^{\prime})\) , every closed subset of a normal space is a normal subspace; but this is not always the case for an arbitrary subset of a normal space.

For example, a completely regular space which is not normal is homeomorphic to a subspace of a compact space (§ 1, no. 5, Proposition 3), and the latter is normal.

Finally we record that the product of two normal spaces is not necessarily normal (see Exercise 9 and § 5, Exercise 16).

